\chapter{Friedmann Expansion, Radiation, and Dark Matter}

\section{Review of the Basic Cosmological Equation}

On scales larger than clusters of galaxies, the universe is taken to be homogeneous. Choose coordinates so that we sit at \(x=0\), and place a representative galaxy, or simply an imaginary particle moving with the Hubble flow, at \(x=1\). Its physical distance from us is then, by definition, the scale factor \(a(t)\). A point at \(x=2\) is at distance \(2a(t)\), and so on.

\begin{figure}[ht]
\centering
\includegraphics[width=0.78\textwidth]{lecture_03_figure_02.png}
\caption{Homogeneous sphere centered at \(x=0\), a comoving galaxy at \(x=1\), and the Newtonian energy balance for \(a(t)\).}
\end{figure}

Only the matter inside the sphere of radius \(a(t)\) contributes to the gravitational pull on the particle at \(x=1\). Give that particle unit mass. Its kinetic energy is
\begin{equation}
\frac{1}{2}\dot a^2.
\end{equation}
The mass enclosed by the sphere is
\begin{equation}
M_{\mathrm{enc}}=\frac{4\pi}{3}a^3\rho,
\end{equation}
so the gravitational potential energy is
\begin{equation}
-\frac{G M_{\mathrm{enc}}}{a}
=
-\frac{4\pi G}{3}a^2\rho.
\end{equation}
\editorialnote{At \lecturetimestamp{00:05:30} the board clearly shows the kinetic term and the density-dependent potential term. The full coefficient written here follows from the spoken derivation.}

Energy conservation gives
\begin{equation}
\frac{1}{2}\dot a^2-\frac{4\pi G}{3}a^2\rho=\text{const}.
\end{equation}
The constant is then written as \(-k\), with the usual warning that the standard normalization really fixes that notation only after multiplying the equation by \(2\). In the standard form,
\begin{equation}
\dot a^2=\frac{8\pi G}{3}a^2\rho-k,
\end{equation}
and dividing by \(a^2\) gives
\begin{equation}
\left(\frac{\dot a}{a}\right)^2=\frac{8\pi G}{3}\rho-\frac{k}{a^2}.
\end{equation}
This is the Friedmann equation in the form used throughout the discussion.

\begin{figure}[ht]
\centering
\includegraphics[width=0.56\textwidth]{lecture_03_figure_03.png}
\caption{Curvature contribution \(-k/a^2\) isolated in the Friedmann equation.}
\end{figure}

The curvature term is empirically very close to zero, close enough that at this level it is a very good approximation simply to neglect it. It does correspond to the curvature of space, but observationally it is too small here to determine even its sign.

The same equation also comes from general relativity. In schematic form, the time-time Einstein equation is
\begin{equation}
R_{00}-\frac{1}{2}g_{00}R \propto T_{00}.
\end{equation}
\editorialnote{The discussion gives only the tensor structure here, not the exact overall numerical coefficient.}
With an expanding homogeneous metric, the geometric side produces the combination of \(\dot a/a\) and \(k/a^2\), while the matter side produces the energy density. At this stage, the Newtonian argument and the Einstein equation are saying the same thing.

That also changes the meaning of \(\rho\). From here on, it is better to call \(\rho\) the energy density. For slowly moving matter in the comoving frame, that is just the rest-energy density.

\section{Matter Domination}

Now assume that the energy density is carried by matter at rest in the comoving frame. Take a cube one coordinate unit on a side. Its physical side lengths are \(a(t)\), so its physical volume is \(a^3(t)\). If the mass-energy inside that comoving cube is \(M\), and matter moves with the flow rather than crossing the boundary, then
\begin{equation}
\rho_m=\frac{M}{a^3}.
\end{equation}

This is the matter-dominated case. In this idealization the particles are at rest in the box, so the temperature is taken to be zero and the pressure is taken to be zero. Pressure would require particles crossing the boundary of the box.

With \(k\approx 0\), the equation becomes
\begin{equation}
\left(\frac{\dot a}{a}\right)^2=\frac{8\pi G}{3}\frac{M}{a^3}.
\end{equation}
Try a power law,
\begin{equation}
a(t)=Ct^p.
\end{equation}
Then
\begin{equation}
\frac{\dot a}{a}=\frac{p}{t},
\end{equation}
so the left-hand side scales like \(t^{-2}\), while the right-hand side scales like \(t^{-3p}\). Matching powers gives
\begin{equation}
3p=2,
\qquad
p=\frac{2}{3}.
\end{equation}
Therefore
\begin{equation}
a(t)\propto t^{2/3}.
\end{equation}
Keeping the coefficient gives
\begin{equation}
C^3=6\pi G M.
\end{equation}

The point of the solution is the power law. A matter-dominated universe expands like the two-thirds power of time, not linearly. This law describes a long intermediate era, from roughly \(10^5\) years after the Big Bang to a few billion years later. The present universe is not literally matter-dominated, because dark energy now dominates, but ordinary matter and a large stretch of cosmological history are described this way.

\section{Radiation Domination}

Earlier than that, the energy density was dominated by radiation. The reason is simple. The universe contains an enormous number of photons, roughly \(10^9\) photons for every proton, and when the universe is compressed those photons become more energetic.

A photon has no rest mass, but it does have energy. Its energy is
\begin{equation}
E_\gamma=h\nu=\frac{hc}{\lambda}.
\end{equation}
The crucial assumption is that as the universe expands, photon wavelengths stretch with it:
\begin{equation}
\lambda\propto a.
\end{equation}
A wave in a box gives the right picture. If the box grows, the wavelength of the mode grows with the box. The same idea is used here for the expanding universe.

Since \(\lambda\propto a\), the energy of each photon scales as
\begin{equation}
E_\gamma\propto a^{-1}.
\end{equation}
Take again a comoving unit box. Photons pass in and out, but on average the number inside stays fixed. Then the total photon energy in the box also scales like \(a^{-1}\). Dividing by the physical volume \(a^3\), the radiation energy density scales as
\begin{equation}
\rho_r\propto a^{-4}.
\end{equation}

That extra factor of \(a^{-1}\) is the whole difference. Matter dilutes because the universe expands. Radiation dilutes for that reason and also because each photon is redshifted.

The Friedmann equation for a radiation-dominated universe is therefore
\begin{equation}
\left(\frac{\dot a}{a}\right)^2=\frac{8\pi G}{3}\frac{M_r}{a^4}.
\end{equation}
\editorialnote{The spoken discussion reuses the symbol \(M\) for this new constant. The subscript \(r\) is added here only to avoid confusion with the matter case.}

Use the same trial form \(a(t)=Ct^p\). The left-hand side again scales like \(t^{-2}\), while the right-hand side scales like \(t^{-4p}\). Matching powers gives
\begin{equation}
4p=2,
\qquad
p=\frac{1}{2}.
\end{equation}
So
\begin{equation}
a(t)\propto t^{1/2}.
\end{equation}
Keeping the coefficient gives
\begin{equation}
C^4=\frac{32\pi G}{3}M_r.
\end{equation}

\begin{classroomqa}
\lecturetimestamp{00:35:41}
\audiencequestion{Why is there a \(G\) in there?}
\lecturerresponse{Because what gravitates is energy density, not rest mass alone. For slowly moving matter the distinction is easy to miss, but photons force it on us: they have no rest mass and still gravitate.}
\end{classroomqa}

A box full of radiation makes the point concrete. Push on a box filled with photons. Once the box accelerates, the wall moving into the photons is struck more strongly than the wall moving away from them, so one has to push harder against the radiation pressure. The box therefore has extra inertia, and by the equivalence principle that extra inertia also means extra gravitational mass. This is the box-of-photons argument used here to say that energy gravitates.

\begin{classroomqa}
\lecturetimestamp{00:38:11}
\audiencequestion{Is either of these equations solvable with an \(a\) that decreases with time?}
\lecturerresponse{Yes. If the curvature term \(k/a^2\) is restored, there are solutions that expand and then contract. That is not what our universe appears to be doing, but the equation allows it.}
\end{classroomqa}

\begin{classroomqa}
\lecturetimestamp{00:38:44}
\audiencequestion{What happened to conservation of energy?}
\lecturerresponse{If the \(k\) term is restored and the equation is multiplied back by \(a^2\), the cosmological equation is precisely the conservation law. Matter or radiation on one side is only part of the bookkeeping. The expansion contributes the kinetic and gravitational part, and the total balance is conserved.}
\end{classroomqa}

\section{Matter and Radiation Together}

Real cosmology contains both matter and radiation. Their densities scale differently:
\begin{equation}
\rho_m\propto a^{-3},
\qquad
\rho_r\propto a^{-4}.
\end{equation}
So at small \(a\), radiation dominates. At large \(a\), matter dominates. Running backward, the early universe must therefore have been radiation-dominated.

\begin{figure}[ht]
\centering
\includegraphics[width=0.68\textwidth]{lecture_03_figure_05.png}
\caption{Axes for comparing matter and radiation energy density as functions of the scale factor.}
\end{figure}

\editorialnote{At \lecturetimestamp{00:44:45} the board shows only the axes at the start of the comparison plot. The matter and radiation curves discussed immediately afterward are described in the text rather than redrawn here.}

The matter curve falls like \(a^{-3}\). The radiation curve falls like \(a^{-4}\). Both diverge as \(a\to 0\), but the radiation curve rises faster there, so the two curves cross at an equality scale factor. A rough estimate places this transition after roughly \(10^4\) years, perhaps a little more.

The same point can be expressed in the plot of \(a(t)\) against \(t\). Early on, the expansion behaves like \(t^{1/2}\). Later it behaves like \(t^{2/3}\). The change is not a jump; it is a continuous crossover through a relatively narrow range of time.

\begin{classroomqa}
\lecturetimestamp{00:44:04}
\audiencequestion{As the universe ages, does the power \(p\) go continuously from \(1/2\) to \(2/3\)?}
\lecturerresponse{Yes. It does not jump. The change is continuous, though it happens over a fairly narrow range of time.}
\end{classroomqa}

\section{Temperature, Recombination, and Last Scattering}

The next quantity is temperature. If the energy of a typical photon scales like \(a^{-1}\), then so does the temperature of the radiation bath:
\begin{equation}
T\propto a^{-1}.
\end{equation}
If \(a\) doubles, the temperature is cut in half.

Take the present photon temperature to be about \(3\,\mathrm K\), and the ionizing temperature to be about \(3000\,\mathrm K\), roughly the temperature of the surface of the sun. Then
\begin{equation}
\frac{a_{\mathrm{today}}}{a_{\mathrm{ion}}}\sim 10^3.
\end{equation}
Using the late-time matter-dominated law \(a\propto t^{2/3}\),
\begin{equation}
\left(\frac{t_{\mathrm{today}}}{t_{\mathrm{ion}}}\right)^{2/3}\sim 10^3,
\end{equation}
so
\begin{equation}
\frac{t_{\mathrm{today}}}{t_{\mathrm{ion}}}\sim 1000^{3/2}\sim 3\times 10^4.
\end{equation}
Taking the age of the universe, at the level of this estimate, to be \(10^{10}\) years gives
\begin{equation}
t_{\mathrm{ion}}\sim 3\times 10^5\ \mathrm{years}.
\end{equation}

This is the epoch when hydrogen and helium stop being fully ionized. Earlier than that, the universe is a hot plasma. Why is a plasma opaque? Because free charges respond strongly to electromagnetic waves. The electric field drives the charges back and forth, the medium behaves like a conductor, and conductors are opaque. That is the relevant analogy with the surface of the sun.

Once the universe cools enough for electrons to recombine with nuclei, ordinary neutral atoms form and the universe becomes much more transparent. The transition is not infinitely sharp. It is a fuzzy place, or a fuzzy time interval, but it is sharp enough for cosmology that it is called the surface of last scattering.

Looking outward means looking backward in time. Nearby galaxies show the recent past. More distant galaxies show earlier times. Far enough out, every line of sight reaches the epoch when the universe was still ionized and opaque. Electromagnetic radiation cannot show anything earlier than that.

Other messengers could in principle probe earlier times. Neutrinos are electrically neutral and would decouple earlier. Gravitons, if one could ever observe them, would come from still earlier epochs. But the same weakness of interaction that would let them pass through the early plasma also makes them extremely hard to detect.

If the universe was once as hot as the surface of the sun in every direction, why is the sky not blindingly bright? Because the last-scattering surface is receding, and the photons arriving from it have been redshifted by roughly the same factor of \(10^3\) that separates \(3000\,\mathrm K\) from \(3\,\mathrm K\). What reaches us is microwave radiation, not visible sunlight.

\begin{classroomqa}
\lecturetimestamp{01:03:51}
\audiencequestion{So our telescopes cannot look earlier than about \(300{,}000\) years?}
\lecturerresponse{Electromagnetic telescopes cannot. Before that time the universe was ionized and opaque. Other messengers, at least in principle, could probe earlier epochs.}
\end{classroomqa}

\begin{classroomqa}
\lecturetimestamp{01:08:43}
\audiencequestion{When you talk about the temperature of a photon, you mean the temperature of a black body with that peak radiation. But that is not linear in wavelength, right?}
\lecturerresponse{Right. The temperature is essentially the energy of a typical photon, so for a photon it scales inversely with wavelength. The \(T^4\) law refers to the energy density, not the energy of a single photon. It combines the energy per photon with the number of photons per unit volume, and the number density scales like \(T^3\).}
\end{classroomqa}

Thus
\begin{equation}
n_\gamma\propto T^3,
\qquad
\rho_\gamma\propto T^4.
\end{equation}

\section{Dark Matter from Rotation Curves}

Dark matter enters through galactic rotation curves. It was not inferred from the large-scale Friedmann equation, but from astronomy on galactic scales.

Start with the simplest Newtonian problem: a light body of mass \(m\) orbiting a central mass \(M\) at radius \(r\). The gravitational force is
\begin{equation}
\frac{GMm}{r^2},
\end{equation}
and the centripetal acceleration is \(v^2/r\). Newton's law gives
\begin{equation}
m\frac{v^2}{r}=\frac{GMm}{r^2}.
\end{equation}
Cancel \(m\), multiply by \(r\), and obtain
\begin{equation}
v(r)=\sqrt{\frac{GM}{r}}.
\end{equation}
So if the gravitating mass is concentrated near the center, the orbital speed falls like \(r^{-1/2}\).

That is what one would expect if mass followed light. Most of the visible matter in a galaxy lies near the core, so outer stars would then behave much like planets in the outer solar system. Astronomers infer the orbital speeds by measuring Doppler-shift differences across a nearby galaxy, comparing the side moving toward us with the side moving away after subtracting the overall recessional motion of the galaxy as a whole. The measured orbital speeds stay roughly constant out to large radii.

So replace the central mass \(M\) by the enclosed mass \(M(r)\):
\begin{equation}
\frac{v^2}{r}=\frac{G M(r)}{r^2}.
\end{equation}
If \(v(r)\) is approximately constant, then
\begin{equation}
M(r)\propto r.
\end{equation}
In other words, the gravitating mass keeps increasing with radius even where the visible light does not.

Now assume, as the discussion does for simplicity, that the mass distribution is roughly spherical. Then
\begin{equation}
M(R)=\int_0^R 4\pi r^2\rho(r)\,dr.
\end{equation}
If \(M(R)\propto R\), differentiating with respect to \(R\) gives
\begin{equation}
4\pi R^2\rho(R)\propto 1,
\end{equation}
and therefore
\begin{equation}
\rho(r)\propto \frac{1}{r^2}.
\end{equation}

\begin{classroomqa}
\lecturetimestamp{01:28:18}
\audiencequestion{It looks as if the density should go like \(1/r^2\). Is there a reason for that?}
\lecturerresponse{From the observed flat rotation curve one infers \(M(R)\propto R\), and differentiating then gives \(\rho(r)\propto 1/r^2\). Beyond that, no simple first-principles explanation is claimed here; detailed simulations do much of the work.}
\end{classroomqa}

The stable conclusion is that a galaxy contains much more gravitating mass than luminous matter alone would suggest, of order several times larger and roughly an order of magnitude larger in the estimate used here. The visible galaxy sits inside a much larger dark halo.

What keeps ordinary gas from staying in such a diffuse state? Dissipation. A gas cloud shrinks because collisions, friction-like processes, and radiation allow it to lose energy. Dark matter evidently does not radiate appreciably and is nearly collisionless. Gravity can pull it inward and convert potential energy into kinetic energy, but without efficient dissipation it does not collapse into the compact luminous structures made by ordinary matter. What keeps the halo extended is kinetic energy.

Several alternatives are then discussed cautiously. Neutrinos are said to be too fast to behave this way on galactic scales. A halo made from ordinary hydrogen would badly upset early-universe nucleosynthesis. Black holes made from ordinary matter are also mentioned as problematic, because they would imply too much ordinary matter in the early universe and would raise lensing issues; the discussion does not claim more than that on the spot.

The halo is not described as a thin shell or a torus. It is an approximately spherical, non-uniform ball, denser toward the center and falling gradually outward.

\begin{classroomqa}
\lecturetimestamp{01:53:12}
\audiencequestion{Are you talking about the Bullet Cluster?}
\lecturerresponse{Yes. The X-ray gas shows the ordinary matter, while lensing shows the total gravitating mass. In the collision those two distributions separate. That is strong evidence for real dark matter rather than a mere modification of gravity tied only to visible matter.}
\end{classroomqa}

That is the final observational point. Very little is known in detail about what dark matter is, but nothing here requires a radical new principle. The conservative possibility is simply that nature contains additional weakly interacting particles that have not yet been identified.
